\chapter{Quotient rings and the Chinese remainder theorem}\label{chap:M07.09}

% Topic grouping: Descent and isomorphisms
\KEntry{K07.09.quotient_construction}{Quotient operations}
\begin{theorem}\label{thm:quotient-ring-construction}
If $I$ is a two-sided ideal of a ring $R$, its additive cosets form a
ring $R/I$ with $[a]+[b]=[a+b]$, $[a][b]=[ab]$, and $1=[1]$.
Conversely, these coset operations on an additive subgroup are
well-defined only if that subgroup is two-sided.
\end{theorem}
\begin{proof}
The additive quotient is an abelian group. If $a'=a+i$ and $b'=b+j$,
then $a'b'-ab=aj+ib+ij\in I$, proving representative independence of
multiplication; addition follows from additive-subgroup closure.
Associativity, distributivity and the identity follow by taking the
class of each corresponding identity in $R$. Conversely, if $[i]=[0]$
for $i$ in the subgroup, well-defined multiplication gives
$[ri]=[r][0]=[0]$ and $[ir]=[0][r]=[0]$ for all $r$. These are exactly
the two absorption conditions.
\end{proof}

\KEntry{K07.09.quotient_projection}{The natural projection}
\begin{proposition}\label{prop:quotient-projection}
The map $\pi:R\to R/I$, $\pi(a)=[a]$, is a surjective unital ring
homomorphism with kernel $I$. Moreover
$[a]=[b]$ if and only if $a-b\in I$.
\end{proposition}
\begin{proof}
Coset equality has the stated additive meaning. In particular $[a]=0$
exactly when $a\in I$. The operation formulas show that $\pi$ preserves
addition, multiplication and $1$, and every coset has a representative.
\end{proof}
This also completes the assertion that every two-sided ideal is a kernel.

\KEntry{K07.09.quotient_trivial}{Trivial and zero-ideal quotients}
$R/I$ is trivial if and only if $I=R$: triviality says $[1]=[0]$,
equivalently $1\in I$, equivalently $I=R$.
The map $R/(0)\to R$ sending $[a]$ to $a$ is well-defined because
$[a]=[b]$ now means $a=b$; its inverse is the natural projection.
Both preserve operations by the quotient formulas. Thus properness is
unnecessary for constructing a quotient, but precisely guarantees that
the quotient is nontrivial.

\KEntry{K07.09.quotient_lift}{Universal factorization}
\begin{theorem}\label{thm:ring-quotient-universal}
For a unital ring homomorphism $f:R\to S$, there is a homomorphism
$\bar f:R/I\to S$ with $\bar f\pi=f$ if and only if $I\subseteq\ker f$.
When it exists it is unique and satisfies $\bar f([a])=f(a)$.
\end{theorem}
\begin{proof}
If a factorization exists, $i\in I$ gives $f(i)=\bar f(0)=0$.
Conversely, $a-b\in I$ implies $f(a)-f(b)=0$, so the proposed formula
is well-defined. Applying $f$ to the quotient operation formulas shows
it is a unital homomorphism. Surjectivity of $\pi$ forces uniqueness.
\end{proof}

\KEntry{K07.09.first_isomorphism}{The first isomorphism theorem}
\begin{theorem}\label{thm:ring-first-isomorphism}
A unital ring homomorphism $f:R\to S$ induces an isomorphism
$R/\ker f\cong\im f$, taking $[a]$ to $f(a)$. If $f$ is onto,
the target is $S$.
\end{theorem}
\begin{proof}
The preceding theorem constructs this homomorphism with codomain
restricted to $\im f$. It is surjective by definition of image.
Its zero class consists of the $a$ with $f(a)=0$, precisely $\ker f$,
so its kernel as a quotient map is zero. The zero-kernel test makes it
injective, and the inverse of a bijective homomorphism is a homomorphism.
\end{proof}

\KEntry{K07.09.second_isomorphism}{The second isomorphism theorem}
Let $A$ be a unital subring of $R$ and $I$ a two-sided ideal.
The set $A+I$ is a unital subring: subtraction preserves it, and
$(a+i)(b+j)=ab+(aj+ib+ij)\in A+I$.
The intersection $A\cap I$ is a two-sided ideal of $A$.
The restriction $A\to(A+I)/I$, $a\mapsto[a]$, is onto, since
$[a+i]=[a]$, and has kernel $A\cap I$.
Theorem~\ref{thm:ring-first-isomorphism} therefore gives
\begin{equation}\label{eq:ring-second-iso}
 A/(A\cap I)\ \cong\ (A+I)/I.
\end{equation}
Here $I$ on the right is regarded as an ideal of $A+I$.
For a nonunital subring the same formulas and proof give a nonunital
isomorphism; they do not turn it into a unital subring.

\KEntry{K07.09.third_isomorphism}{The third isomorphism theorem}
For two-sided ideals $I\subseteq J$, the subset
$J/I=\{[j]:j\in J\}$ is an ideal of $R/I$: differences are represented
by differences in $J$, and multiplication by $[r]$ on either side
uses absorption in $J$. The map $R/I\to R/J$, $[a]_I\mapsto[a]_J$,
is well-defined and onto, with kernel $J/I$.
The first isomorphism theorem proves
\[
 (R/I)/(J/I)\ \cong\ R/J.
\]
No commutativity has been used.

\KEntry{K07.09.ideal_correspondence}{The ideal correspondence}
\begin{theorem}\label{thm:ideal-correspondence}
For a surjective unital ring homomorphism $f:R\to S$, taking inverse
image is an inclusion-preserving bijection from two-sided ideals of $S$
to two-sided ideals of $R$ containing $\ker f$. Its inverse is taking
image. In particular ideals of $R/I$ correspond to ideals above $I$.
\end{theorem}
\begin{proof}
Inverse images preserve subtraction and both absorptions; they contain
$\ker f$. If $J$ is an ideal above the kernel, then $f(J)$ is an ideal:
to multiply $f(j)$ by $s\in S$, choose $r$ with $f(r)=s$ and use
$rj,jr\in J$. Further $f^{-1}(f(J))=J$: an equality $f(x)=f(j)$
gives $x-j\in\ker f\subseteq J$. Surjectivity gives
$f(f^{-1}(H))=H$ for an ideal $H$ of $S$.
Both operations preserve inclusion, proving the claimed ordered bijection.
\end{proof}
Without surjectivity an image need not be an ideal of the target:
the image of $\Z\hookrightarrow\Q$ is not an ideal of $\Q$.
The ideal generated by an image is a different operation.

% Topic grouping: Prime and maximal ideals
Until the end of this chapter rings are commutative.
\KEntry{K07.09.prime_ideal}{Prime ideals}
A prime ideal $P$ is proper and satisfies
$ab\in P\Rightarrow a\in P$ or $b\in P$.
Equivalently, for ideals $I,J$, $IJ\subseteq P$ implies
$I\subseteq P$ or $J\subseteq P$. To prove this direction choose
$a\in I\setminus P$, if possible; for every $b\in J$ the product lies
in $P$, so $b\in P$. Conversely apply the ideal version to $(a),(b)$.
Their product is $(ab)$, since every summand $(ra)(sb)$ is a multiple
of $ab$, and every multiple $t ab=(ta)b$ belongs to the product.

\KEntry{K07.09.maximal_ideal}{Maximal ideals}
A maximal ideal is a proper ideal $M$ with no ideal strictly between
$M$ and $R$. In a noncommutative ring this must be distinguished from
a maximal \emph{left} ideal. The quotient-field statements below are
commutative statements.

\KEntry{K07.09.prime_quotient}{Domains as prime quotients}
\begin{proposition}\label{prop:prime-quotient}
$P$ is prime if and only if $R/P$ is an integral domain.
\end{proposition}
\begin{proof}
Properness is exactly nontriviality of $R/P$. Also
$[a][b]=0$ is equivalent to $ab\in P$, and $[a]=0$ is equivalent
to $a\in P$. Translating the zero-product condition gives precisely
the prime-ideal condition.
\end{proof}
In particular $R$ is a domain if and only if $(0)$ is prime.

\KEntry{K07.09.maximal_quotient}{Fields as maximal quotients}
\begin{proposition}\label{prop:maximal-quotient}
$M$ is maximal if and only if $R/M$ is a field.
\end{proposition}
\begin{proof}
The ideal correspondence identifies ideals of $R/M$ with those between
$M$ and $R$. If $M$ is proper, the quotient is nontrivial.
Consequently maximality is equivalent to its having only zero and whole
ideals, which is the field criterion
of Proposition~\ref{prop:field-ideal-criterion}.
\end{proof}

\KEntry{K07.09.maximal_prime_nonconverse}{Prime does not mean maximal}
A field is a domain, so the two quotient criteria prove that maximal
ideals are prime. The zero ideal of $\Z$ is prime but not maximal,
since $\Z$ is a domain and not a field. A nonzero example is $(X)$ in
$\Z[X]$: evaluation at $0$ is onto $\Z$ with kernel $(X)$, so
$\Z[X]/(X)\cong\Z$ is a domain and not a field.
The kernel assertion follows directly from coefficients:
$f=f(0)+Xg$, where $g$ shifts the positive-index coefficients down
by one, so $f(0)=0$ exactly when $f$ is a multiple of $X$.
Equivalently $(X)\subsetneq(2,X)\subsetneq\Z[X]$.
The first inclusion is strict because the constant $2$ is not in $(X)$;
the second because a member $2f+Xg$ has even constant coefficient.

\KEntry{K07.09.maximal_existence}{Existence of maximal ideals}
\begin{theorem}\label{thm:maximal-ideal-existence}
Every proper ideal is contained in a maximal ideal.
\end{theorem}
\begin{proof}
Order the proper ideals above a fixed proper $I$ by inclusion.
An empty chain has upper bound $I$. A nonempty chain has its union
as an ideal upper bound. That union is proper: otherwise it contains
$1$, which belongs to a single member, contradicting its properness.
Zorn's lemma gives a maximal member. Any larger proper ideal would still
contain $I$, so this member is maximal among all proper ideals.
\end{proof}
This uses choice and is not an algorithm to find a maximal ideal.

\KEntry{K07.09.contraction}{Inverse images}
If $f:R\to S$ is unital and $P\subset S$ prime, $f^{-1}(P)$ is proper
because $f(1)=1\notin P$. If $ab$ maps into $P$, one of $f(a),f(b)$
does, proving primality. If $f$ is onto and $M$ is maximal, the ideal
correspondence or the isomorphism
$R/f^{-1}(M)\cong S/M$ proves maximality of its inverse image.
Surjectivity cannot be omitted: under $\Z\hookrightarrow\Q$, the
maximal ideal $(0)$ of $\Q$ contracts to the nonmaximal $(0)$ of $\Z$.
For nonunital maps even properness of a prime inverse image can fail.

% Topic grouping: Simultaneous congruences
\KEntry{K07.09.comaximal}{Comaximality}
Ideals $I,J$ are comaximal when $I+J=R$, equivalently when
$i+j=1$ for some $i\in I,j\in J$.
Distinct maximal ideals are comaximal: if $J\not\subseteq I$, maximality
of $I$ forces $I+J=R$; if $J\subseteq I$, maximality of $J$ would give
$J=I$, since $I$ is proper. For principal ideals this is the existence
of a Bezout identity, a stronger condition than merely having gcd $1$
in an arbitrary domain.

\KEntry{K07.09.product_intersection}{Comaximal products}
\begin{proposition}\label{prop:comaximal-product-intersection}
Comaximal ideals satisfy $IJ=I\cap J$. For a finite pairwise comaximal
family, its product equals its intersection.
\end{proposition}
\begin{proof}
The product is always contained in the intersection. If $x\in I\cap J$
and $i+j=1$ as above, then $x=xi+xj\in IJ$, using commutativity.
For the finite assertion, if $I_1+I_j=R$, choose
$a_j\in I_1,b_j\in I_j$ with $a_j+b_j=1$.
Expanding $\prod_{j=2}^m(a_j+b_j)=1$ gives
$1\in I_1+I_2\cdots I_m$: all terms except $\prod b_j$ have a factor
in $I_1$. Induction now applies the binary result to $I_1$ and the
product of the others.
\end{proof}

\KEntry{K07.09.crt_binary}{The binary Chinese remainder theorem}
\begin{theorem}\label{thm:crt-binary}
For comaximal ideals $I,J$, the map
$r\mapsto([r]_I,[r]_J)$ induces
\[
 R/(IJ)\ \cong\ R/I\times R/J.
\]
If $i+j=1$ with $i\in I,j\in J$, then $aj+bi$ represents the inverse
image of $([a]_I,[b]_J)$.
\end{theorem}
\begin{proof}
The displayed map is a unital homomorphism by coordinatewise
operations, with kernel $I\cap J$. Since $j\equiv1\pmod I$ and
$i\equiv1\pmod J$, its value at $aj+bi$ is the required pair.
It is therefore onto. Apply the first isomorphism theorem and
Proposition~\ref{prop:comaximal-product-intersection}.
\end{proof}
Without comaximality surjectivity can fail: with $I=J=(2)$ in $\Z$,
the image consists only of equal pairs of residues.

\KEntry{K07.09.crt_finite_integer}{Finite and integer CRT}
\begin{theorem}\label{thm:crt-finite}
For a finite pairwise comaximal family $I_1,\ldots,I_m$,
\[
 R/\bigcap_i I_i\ \cong\prod_i R/I_i,\qquad
 \bigcap_iI_i=\prod_i I_i.
\]
\end{theorem}
\begin{proof}
For $m=0$ both sides are trivial: the intersection and empty ideal
product are $R$, and the empty ring product has one element.
For $m=1$ the assertion is the identity map. For larger $m$, the proof
of the preceding proposition makes $I_1$ comaximal with
$I_2\cdots I_m$. Apply the binary CRT, then the induction isomorphism
to the second factor. Every resulting coordinate map is reduction
modulo the corresponding $I_i$, so this is the claimed natural map.
\end{proof}
For pairwise coprime positive integers $n_i$, the integer Bezout identity
makes $(n_i)+(n_j)=\Z$. Their ideal product is $(\prod_i n_i)$.
Thus any residues $a_i$ have a simultaneous solution, unique modulo
$N=\prod_i n_i$. More explicitly let $N_i=N/n_i$.
Since $N_i,n_i$ are coprime, choose $t_i$ with
$t_iN_i\equiv1\pmod{n_i}$. Then
$\sum_i a_it_iN_i$ has every prescribed residue: all other summands
vanish in that coordinate. This is the D08 integer CRT as a specialization
of one ring theorem, not a separate unrelated construction.

% Topic grouping: A prime-separation lemma
\begin{lemma}\label{lem:nilpotent-prime-separation}
For a commutative ring, an element belongs to every prime ideal if
and only if it is nilpotent.
\end{lemma}
\begin{proof}
If $x^n=0$, repeated primality forces $x$ into any prime ideal.
Conversely suppose no power of $x$ is zero. The multiplicative set
$S=\{1,x,x^2,\ldots\}$ avoids $0$. Zorn's lemma gives an ideal $P$
maximal among ideals disjoint from $S$: the zero ideal belongs to this
partially ordered set and bounds its empty chain; nonempty chain unions remain disjoint,
since an element in the union lies in one member.
If $ab\in P$ but $a,b\notin P$, both $P+(a)$ and $P+(b)$ meet $S$.
Write $x^r=p+ua$, $x^s=q+vb$ with $p,q\in P$.
Multiplication gives $x^{r+s}\in P$, a contradiction. Thus $P$ is
prime, and $x\notin P$ because $x\in S$. This proves the contrapositive.
The zero ring causes no exception: it has no prime ideals and its
single element is nilpotent.
\end{proof}
